%! TEX root = master.tex
\chapter{5}{Real and Complex Differentiability}

Complex differentiability looks like the real definition with $x$ replaced by
$z$, but it is far more rigid. Here are two places where the theories part
ways.

\section{One derivative gives all of them}\label{sec:all-derivatives}

On $\R$, a function can be differentiable once and no more. On $\C$, an
analytic function has derivatives of every order, all analytic
(\cref{thm:cif-derivatives}).

\begin{exercise}\label{ex:x-five-thirds}
  Show that $f(x) = x^{5/3}$ is differentiable on $\R$, but $f'$ is not
  differentiable at $0$.
\end{exercise}

\begin{solution}
  For $x \neq 0$, $f'(x) = \frac53 x^{2/3}$. At $0$,
  \[
    f'(0) = \lim_{h \to 0} \frac{h^{5/3}}{h} = \lim_{h \to 0} h^{2/3} = 0,
  \]
  which agrees with $\frac53 x^{2/3}$ at $x = 0$. But
  \[
    \frac{f'(h) - f'(0)}{h} = \frac53 h^{-1/3}
  \]
  is unbounded as $h \to 0$, so $f''(0)$ does not exist.
\end{solution}

\begin{remark}
  No analytic function behaves like this: once $f'$ exists on an open set,
  so do $f''$, $f'''$, and the rest.
\end{remark}

\section{No mean value theorem}\label{sec:no-mvt}

For a real function differentiable on $[a, b]$, some $c$ between $a$ and $b$
has $f'(c) = \frac{f(b) - f(a)}{b - a}$. The special case $f(a) = f(b)$ is
Rolle's theorem. Neither survives in $\C$.

\begin{exercise}\label{ex:z-cubed-rolle}
  Let $g(z) = z^3$ and $\omega = -\frac12 + \frac{\sqrt3}{2} i$. Show that
  $g(1) = g(\omega)$, but $g'$ does not vanish anywhere on the segment from
  $1$ to $\omega$.
\end{exercise}

\begin{solution}
  $\omega = e^{2\pi i/3}$, so $g(\omega) = e^{2\pi i} = 1 = g(1)$. If the mean
  value theorem held, some $c$ on the segment would have
  $g'(c) = 3c^2 = 0$, that is, $c = 0$. But the segment consists of the points
  $1 + t(\omega - 1)$, $0 \leq t \leq 1$, with imaginary part
  $\frac{\sqrt3}{2} t$, which is zero only at $t = 0$, the point $1$. So $0$
  is not on the segment.
\end{solution}

\begin{remark}
  The exponential gives another example: $e^0 = e^{2\pi i}$, yet $e^z$ is
  never zero.
\end{remark}
